\chapter{Non-Parametric Econometric Validation: NPIV Estimation \& KRR Adaptive Kernels}
\label{ch:non_parametric_validation}

A core empirical hurdle in simulating financial panic lies in validating the outcomes. How do we scientifically prove that a simulated 5\% flash crash accurately models the microstructural mechanics of a real-world event? Standard econometric models fail because they rely heavily on rigid parametric assumptions (e.g., normally distributed returns) that collapse entirely during black-swan events. 

Furthermore, as the simulator scales from 400 generative teacher agents to 1,000,000 distilled policy instances (Chapter \ref{ch:distilled_swarm_scaling}), the validation framework must rigorously distinguish genuine emergent macroeconomic behavior from reinforcement learning policy exploitation and simulator artifacts. This chapter details how AMPS relies on Non-Parametric Instrumental Variables (NPIV) and Kernel Ridge Regression (KRR) with diagonal adaptive kernels to bridge distilled policy swarms with rigorous empirical econometrics.

\section{The Limitations of Parametric Microstructure Models}

Traditional mathematical finance is dominated by parametric models like the Black-Scholes options pricing framework or Generalized Autoregressive Conditional Heteroskedasticity (GARCH) volatility models. These models force empirical data into predefined closed-form equations parameterized by fixed statistical moments. 

During systemic liquidity panic, asset returns are fundamentally non-Gaussian; they exhibit severe leptokurtosis (fat tails), negative skewness, and volatility clustering. Price paths are not continuous diffusions, but jump-diffusion processes punctuated by discrete liquidity voiding. If AMPS attempts to tune its student policies against parametric guardrails, the simulation will artificially constrain the panic, failing to replicate the true, chaotic breakdown of liquidity.

\section{Non-Parametric Instrumental Variables (NPIV) Under Misspecification}

A profound empirical challenge in market microstructure is \textbf{Endogeneity}. When an algorithmic market maker widens its bid-ask spread, does it do so because the order flow has become toxic, or did the spread widening itself induce toxic momentum selling? This bidirectional feedback loop makes it impossible to directly estimate structural policy decision functions using classical ordinary least squares (OLS) regression.

To resolve this endogeneity under model misspecification, AMPS implements \textbf{Non-Parametric Instrumental Variables (NPIV)} estimation, drawing on the foundational conditional moment restriction framework of Newey and Powell \cite{newey2003instrumental} and its modern reproducing kernel Hilbert space (RKHS) formulation by Singh, Sahani, and Gretton \cite{singh2019kernel}. 

\subsection{Structural Equation and Exclusion Restriction}

Consider the non-parametric structural relationship modeling the agent's action $Y_t$ (e.g., spread adjustment or aggressive market sell volume) as an unknown non-linear function $h_0(X_t)$ of endogenous market state $X_t$ (such as Order Flow Imbalance and depth volatility), subject to structural unobserved noise $\epsilon_t$:

\begin{equation}
    Y_t = h_0(X_t) + \epsilon_t, \quad \mathbb{E}[\epsilon_t | X_t] \ne 0
\end{equation}

To identify $h_0$, we introduce an exogenous instrumental variable $Z_t$. In AMPS, the \textbf{Randomly Assigned API Rate Limit Throttle} (specifically, simulated HTTP 429 delays, $Z_{\text{HTTP429}}$) serves as an ideal instrument:

\begin{enumerate}
    \item \textbf{Instrument Relevance:} The assigned HTTP 429 throttling delay directly throttles the execution timing and I/O buffer submission of agent decisions, exogenously perturbing the realized order arrival intensity $\lambda_{\text{arrival}}(t)$.
    \item \textbf{Strict Exclusion Restriction:} Simulated HTTP 429 delays are generated by an uncorrelated pseudorandom generator residing outside the economic state loop. They do not enter the agent's utility function, do not alter price-time priority rules in the matching engine, and share zero structural covariance with unmeasured microstructural error terms $\epsilon_{\text{LOB}}$:
    \begin{equation}
        \text{Cov}(Z_{\text{HTTP429}}, \, \epsilon_{\text{LOB}}) = 0
    \end{equation}
\end{enumerate}

\subsection{Spectral Regularization for Ill-Posed Inverses}

Estimating $h_0$ via conditional moment equations $\mathbb{E}[Y_t - h_0(X_t) | Z_t] = 0$ is an ill-posed inverse problem \cite{newey2003instrumental}. The conditional expectation operator $T: L^2(F_X) \to L^2(F_Z)$ is a compact Hilbert-Schmidt operator whose inverse is discontinuous. 

AMPS applies Tikhonov spectral regularization directly within its Numba JIT estimation pipeline, building on kernel instrumental regression methods proven to achieve minimax optimal convergence rates under ill-posed inverse operators \cite{singh2019kernel, meunier2024minimax}:

\begin{equation}
    \hat{h}_\alpha = \arg\min_{h \in \mathcal{H}} \left\{ \frac{1}{N} \sum_{i=1}^N \left( \hat{\mathbb{E}}[Y | Z = z_i] - \hat{T}h(z_i) \right)^2 + \alpha_N \|h\|_{\mathcal{H}}^2 \right\}
\end{equation}

where $\alpha_N > 0$ is a decaying regularization parameter. This spectral dampening prevents high-frequency numerical instability, isolating the true causal price impact generated specifically by the macroeconomic information shock from endogenous market feedback.

\section{Kernel Ridge Regression (KRR) with Diagonal Adaptive Kernels}

To track non-stationary regime shifts across multi-patch stress scenarios without imposing rigid distribution assumptions, AMPS utilizes \textbf{Kernel Ridge Regression (KRR) with Diagonal Adaptive Feature Kernels}, rooted in the classical RKHS regularization principles of Sch{\"o}lkopf and Smola \cite{scholkopf2002learning} and vector-valued operator weighting formalized by Micchelli and Pontil \cite{micchelli2005learning} (with adaptive diagonal metric models analyzed in Li and Lin \cite{li2025diagonal}).

Classical KRR projects input states $x \in \mathcal{X}$ into a reproducing kernel Hilbert space $\mathcal{H}_K$ via an isotropic kernel $K(x, x') = \exp(-\|x - x'\|^2 / 2\sigma^2)$ \cite{scholkopf2002learning}. However, during liquidity cascades, feature sensitivities diverge rapidly: Order Flow Imbalance (OFI) sensitivities may expand by orders of magnitude while historical moving average signals become uninformative.

The Diagonal Adaptive Kernel model resolves this by parameterizing an anisotropic diagonal metric tensor $W = \text{diag}(w_1, w_2, \dots, w_d)$ \cite{micchelli2005learning, li2025diagonal}:

\begin{equation}
    K_W(x, x') = \exp\left( -\frac{1}{2} (x - x')^T W (x - x') \right) = \exp\left( -\frac{1}{2} \sum_{j=1}^d w_j (x_j - x'_j)^2 \right)
\end{equation}

The diagonal weights $w_j \ge 0$ correspond to adaptive eigenvalues learned dynamically during evaluation by minimizing localized cross-validation error across rolling replay windows. When a market enters the Lehman Liquidity Black Hole regime (Chapter \ref{ch:multi_patch_harvesting}), the adaptive kernel dynamically inflates the weights corresponding to top-of-book depth voiding and OFI skew while dampening trend features.

By pairing NPIV causal estimation with diagonal adaptive KRR, AMPS establishes an automated econometric auditing harness. This framework validates that the 1,000,000 Tier 2 distilled student agents preserve realistic, human-like bounded rationality under stress without devolving into ungrounded simulation heuristics.
